\chapter{Verification, Validation and Reproducibility}
\label{app:E}

\section{The within-class ordering check}
\label{sec:app-priority-check}

The most consequential verification performed in this study concerned the within-class ordering
of the Urgent-First baseline, because the semantics of that policy determine how its results
must be interpreted (\Cref{sec:urgent-first}).

The baseline inserts orders into a priority queue as priority-tagged items, with urgent orders
at priority \num{0} and normal orders at priority \num{1}. The natural assumption --- and the
one the policy's name invites --- is that equal-priority orders are then served in
first-in-first-out order. The comparison operator of the priority-tagged item type, however,
compares \emph{only} the priority field, so equal-priority items are never ordered relative to
one another and their sequence is determined by the internal structure of the binary heap.

Rather than infer this from the source, it was established by direct observation. An
observational script inserts ten items of equal priority in the order \(1, 2, \dots, 10\) and
records the order in which they are returned:

\begin{center}
\begin{tabularx}{0.9\textwidth}{lX}
\toprule
Insertion order & \(1,\ 2,\ 3,\ 4,\ 5,\ 6,\ 7,\ 8,\ 9,\ 10\) \\
Dequeue order   & \(1,\ 3,\ 7,\ 10,\ 9,\ 6,\ 8,\ 5,\ 2,\ 4\) \\
\midrule
First-in-first-out within class? & No \\
Sorted by order identifier?      & No \\
\bottomrule
\end{tabularx}
\end{center}

A second test interleaves the two classes and confirms that the \emph{class-level} priority is
strictly respected: every urgent order is returned before any normal order. Within each class,
however, the ordering is again neither first-in-first-out nor sorted.

The conclusion carried into \Cref{ch:methodology} and \Cref{ch:discussion} is therefore that
Urgent-First implements strict two-class priority with unspecified within-class ordering. The
script is retained in the repository under \code{thesis\_support/analysis/} and reproduces the
result on demand.

\section{Other verification checks}

\begin{itemize}
  \item \textbf{Utilisation bound.} Stage utilisation is computed as busy worker-minutes over
        available worker-minutes and asserted not to exceed one. A worker cannot be busy longer
        than the operating horizon, so a value above one would indicate an upstream error rather
        than a legitimate result.
  \item \textbf{Termination.} The simulation is bounded by a horizon timeout rather than by a
        completion event, so it terminates regardless of how much work remains. Orders left
        queued at the horizon are unresolved backlog, not a deadlock.
  \item \textbf{Backlog accounting.} Unfinished orders have undefined system time, so every
        comparison against the service threshold evaluates false and they are counted as
        failures. They are additionally reported as a separate count and share.
  \item \textbf{Common random numbers.} Service-time lookups were confirmed to depend only on
        order identifier, stage and scenario seed --- never on dequeue order --- so that all
        three policies encounter identical service times.
  \item \textbf{Queue measurement.} Queue length is accumulated from enqueue and dequeue events
        at the exact simulation times they occur, giving an exact time-weighted average
        independent of event spacing.
\end{itemize}

\section{Validation}

Validation here is internal and structural rather than external and empirical, and the
distinction should not be blurred.

\textbf{What was checked.} The model reproduces behaviours that queueing theory predicts:
waiting times grow sharply as utilisation approaches one; congestion concentrates at the stage
with the highest load ratio; capacity added at a non-constraining stage produces little
improvement; larger workforces reduce backlog monotonically. Relative comparisons behave
sensibly: policies that prioritise a class improve that class's attainment, and the cost model
responds to its inputs as specified.

\textbf{What was not done.} No comparison against a real warehouse was performed. Service-time
distributions were not fitted to observed handling data. Economic parameters were not calibrated
against real labour costs or contractual penalties. No practitioner reviewed the recommendations
for operational plausibility. The model is therefore credible as an instrument for comparing
configurations under stated assumptions, and is \emph{not} established as predictive of any
particular operation.

\section{Reproducibility}

\subsection{What can be reproduced}

\begin{table}[htbp]
\centering
\caption{Reproducibility status of the thesis artefacts.}
\label{tab:repro}
\small
\begin{tabularx}{\textwidth}{lXl}
\toprule
\textbf{Artefact} & \textbf{Source} & \textbf{Status} \\
\midrule
All figures and tables in \Cref{ch:results} & scripts under
  \code{thesis\_support/analysis/} reading persisted run outputs & reproducible \\
Within-class ordering check & \code{verify\_priority\_store.py} & reproducible \\
Capacity-pressure band statistics & \code{verify\_pressure\_bands.py} reading the persisted
  retrospective results & reproducible \\
Demand dataset & seasonal generator with fixed base seed & reproducible \\
FIFO and Urgent-First runs & simulation engine with recorded seeds & reproducible \\
RL-3 runs & requires the trained checkpoint & \textbf{not reproducible} \\
\bottomrule
\end{tabularx}
\end{table}

\subsection{The checkpoint limitation}

As recorded in \Cref{app:C}, the trained network weights matched a repository ignore rule, were
never version-controlled, and were absent from the working tree when this thesis was prepared.
The persisted evaluation outputs are available, and all RL-3 results derive from them, but exact
regeneration of those results requires the trained checkpoint, which was not available in the
environment used for final thesis preparation. The recorded digest and file size document which
artefact produced them; they do not enable it to be recreated. Retraining from the recorded
configuration would yield a different policy, not the same one.

Three points bound the significance of this. First, it affects one of the three compared
policies; the FIFO and Urgent-First results are fully reproducible. Second, the persisted RL-3
outputs are internally consistent with the rest of the run --- the same order data, the same
configurations, the same common random numbers and the same economic parameters --- and were
written by the same pipeline in the same execution. Third, provenance metadata recorded the
checkpoint's size, digest and successful loading into the current architecture.

None of that is a substitute for being able to re-run the evaluation, and the limitation is
reported rather than minimised. Retaining trained model artefacts under version control, or in
an accompanying archive, is the obvious remedy and is recommended for any continuation of this
work.

\subsection{Environment}

\begin{table}[htbp]
\centering
\caption{Environment for the thesis analysis.}
\label{tab:environment}
\small
\begin{tabularx}{0.86\textwidth}{lX}
\toprule
\textbf{Item} & \textbf{Value} \\
\midrule
Source commit        & \code{c42c8e4} \\
Python               & 3.12 \\
Key packages         & SimPy 4.1.1, NumPy 2.3.2, pandas 2.3.1, Matplotlib 3.10.7 \\
Retrospective run    & \code{data/api\_runs/latest/historical/} \\
Forecast run         & \code{data/api\_runs/latest/future/} \\
Retrospective dataset & \num{240000} orders, scenario \code{seasonal\_base} \\
Forecast streams     & generated per replication, scenario
                       \code{future\_december\_standard\_r0..r2} \\
Analysis scripts     & \code{thesis\_support/analysis/} \\
\bottomrule
\end{tabularx}
\end{table}

All analysis scripts written for this thesis are observational: they read persisted outputs and
produce figures and tables. None modifies the simulation, the configuration or the run outputs.
